\subsubsection{Análisis por clase y errores en ISIC2019}
La Figura~\ref{fig:errores-isic2019} presenta a la izquierda la matriz de confusión de B71-limpio y a la derecha precisión, sensibilidad y F1 por diagnóstico. Las filas representan clases reales y las columnas clases predichas; cada fila suma el 100\,\%. La diagonal expresa sensibilidad por clase, no el número de imágenes. Los soportes aparecen junto a las métricas. Este análisis corresponde a los pesos limpios de la Tabla~\ref{tab:clasificacion-isic2019}.

\begin{figure}[!htbp]
\caption{Matriz de confusión y métricas por diagnóstico de B71-limpio en ISIC2019.}\label{fig:errores-isic2019}
\centering
\includesvg[width=\linewidth]{images/organizadas/37_errores_isic2019}
\par\smallskip
{\singlespacing\fontsize{11}{13}\selectfont\raggedright
\noindent\textit{Nota.} Elaboración propia desde los registros canónicos; n=5.067. Porcentajes normalizados por fila. Métricas puntuales, sin intervalos añadidos.\par}
\end{figure}
\FloatBarrier
En ISIC2019, B71-limpio identifica 613 de los 905 melanomas (sensibilidad 0,6773). Confunde 154 melanomas con nevos y 70 con lesiones similares a queratosis benignas. El desempeño también varía entre las clases minoritarias: DF y VASC tienen 47 y 50 imágenes, respectivamente. Esta distribución exige interpretar las métricas junto con su soporte, sin equiparar exactitud global y sensibilidad de cada diagnóstico.

\subsubsection{Discriminación de melanoma en ISIC2019}
La Figura~\ref{fig:curvas-isic2019} muestra las curvas precisión--sensibilidad y ROC de B71-limpio al considerar melanoma frente al resto. Cada punto corresponde a un umbral de probabilidad; las curvas no representan únicamente la predicción por máximo de probabilidad de la matriz anterior. La línea horizontal del panel (a) indica la proporción de melanoma en validación y la diagonal del panel (b) representa ausencia de discriminación. La precisión media (AP) es 0,7354; esta magnitud corresponde a melanoma y no debe confundirse con F1 macro ni con AUC multiclase.

\begin{figure}[!htbp]
\caption{Curvas precisión--sensibilidad y ROC para melanoma de B71-limpio en ISIC2019.}\label{fig:curvas-isic2019}
\centering
\includesvg[width=\linewidth]{images/organizadas/38_curvas_isic2019}
\par\smallskip
{\singlespacing\fontsize{11}{13}\selectfont\raggedright
\noindent\textit{Nota.} Elaboración propia desde probabilidades por imagen. Melanoma: 905 de 5067 imágenes. Curvas empíricas sin suavizado ni intervalos de incertidumbre.\par}
\end{figure}
\FloatBarrier
Las curvas describen discriminación de melanoma en esta partición. Su lectura complementa las métricas a un umbral fijo, pero no demuestra utilidad clínica ni validación externa independiente de HAM10000.

\subsubsection{Calibración y confianza en ISIC2019}
La Figura~\ref{fig:calibracion-isic2019} compara el control público con protocolo B71, B71 original y B71-limpio. La fila superior relaciona confianza media y proporción de aciertos; la inferior muestra cuántas imágenes contribuyen a cada intervalo. Esta disposición permite distinguir una diferencia de calibración de una diferencia en la distribución de confianza. Los intervalos sin imágenes se omiten de la confiabilidad; las cruces señalan intervalos con menos de diez imágenes.

\begin{figure}[!htbp]
\caption{Confiabilidad y distribución de confianza de las configuraciones B71 en ISIC2019.}\label{fig:calibracion-isic2019}
\centering
\includesvg[width=\linewidth]{images/organizadas/39_calibracion_isic2019}
\par\smallskip
{\singlespacing\fontsize{11}{13}\selectfont\raggedright
\noindent\textit{Nota.} Elaboración propia desde los registros de probabilidades. n=5067; 15 intervalos uniformes, ECE ponderado por frecuencia. Misma escala de conteos entre columnas; sin recalibración posterior ni IC. B71 original se distingue de los pesos limpios.\par}
\end{figure}
\FloatBarrier
El ECE de B71-limpio es 0,0498, frente a 0,0706 del control público. Estos valores resumen discrepancias agregadas y no constituyen un contraste estadístico entre configuraciones. Un punto por debajo de la diagonal indica confianza superior a la frecuencia observada de aciertos.

